\documentclass[10pt,a4paper]{amsart}
\usepackage[margin=19mm]{geometry}
\usepackage{amsmath,amssymb,mathtools}
\usepackage[scr=rsfs,cal=euler]{mathalfa}
\usepackage{xfrac}
\usepackage[unicode,hidelinks,bookmarks=false]{hyperref}
\newcommand{\ie}{i.e.\ }
\newcommand{\cf}{cf.\ }
\newcommand{\resp}{respectively\ }
\newcommand{\Subsection}{\subsection}
\providecommand{\nobreakdash}{\nobreak\hbox{-}}
\setlength{\emergencystretch}{2em}
\multlinegap0pt
\newtheorem{lemma}{Lemma}
\newtheorem{theorem}{Theorem}
\title{Thermodynamic Dictionary for Syracuse-like Maps}
\author{Eduardo Santana}
\date{Source: arXiv:2601.03297v8 (CC BY 4.0)}
\begin{document}
\maketitle

\noindent\emph{Source and licence.} Thermodynamic Dictionary for Syracuse-like Maps. Version arXiv:2601.03297v8 (CC BY 4.0), distributed by arXiv under the Creative Commons Attribution 4.0 International licence, http://creativecommons.org/licenses/by/4.0/, as recorded in the arXiv licence metadata for this exact version and shown on the abstract page https://arxiv.org/abs/2601.03297v8. Full licence notice: \texttt{licenses/arxiv-2601.03297-CC-BY-4.0.txt}.\par\smallskip

\noindent\textit{Editorial scope.} Counted unit: the article's theorem finite\_Syracuse (source0.tex lines 465--467). The article's complete original proof of it is delivered with it (source0.tex lines 468--502, one \texttt{proof} environment of 35 lines). No mathematics is written, repaired or re-derived: the statement and the proof are the article's own lines, in the article's own notation, and no retained line is substituted or re-flowed.  Retained uncounted as the source-local context the proof needs: lines 354--356.  Nothing was omitted from any retained span, so the package carries no omission record.  No retained line contains a citation, so the wrapper carries no bibliography and no bibliography entry.  No retained line contains a figure environment, an image-inclusion command or drawing code, so no image is delivered and the package declares no asset dependency.  Editorial formatting only, in addition: an AMS \texttt{amsart} preamble that restates the article's own declarations and macros used by the retained lines, namely the environments \texttt{lemma}, \texttt{theorem} on separate counters, together with the standard packages those lines need.  The retained environments print in this document as Lemma 1, Theorem 1 in order of appearance, whereas the article prints the same items with the numbering of its own full document.  The complete native source of the article as distributed in the arXiv e-print for this exact version is bundled unchanged as \texttt{sources/192218/source0.tex}.
	\begin{lemma}\label{converse}
		Let $f: X \to X$ be a general map. Let us assume that some continuous potential $\phi : X \to \mathbb{N}_{0}$ with respect to $\mathcal{T}$ and with finite pressure does not possess any equilibrium state then we have infinitely many cycles with unbounded set of periods. Conversely, if the set of periods is unbounded, we can customize a continuous potential $\phi : X \to \mathbb{N}_{0}$ with respect to $\mathcal{T}$ and with finite pressure which does not possess any equilibrium state.
	\end{lemma}

\begin{theorem}[Finiteness of cycles for $S_{m,n}$] \label{finite_Syracuse}
	Let $S_{m,n} : \mathbb{N} \to \mathbb{N}$ be a Syracuse map. There exist at most finitely many cycles.
\end{theorem}

\begin{proof}
	We begin by defining a map $\pi : \mathbb{N} \to \Sigma_{2}$ where $\Sigma_{2} = \{0,1\}^{\mathbb{N}}$ and $S : \Sigma_{2} \to \Sigma_{2}$ is the left-sided shift $S(x_{1}, x_{2}, \cdots, x_{k},\cdots) =  (x_{2}, x_{3}, \cdots, x_{k},\cdots)$. With the topology of cylinders and the standard metric, the shift map $S$ is continuous. We define $\pi(x) = (x_{1}, x_{2}, \cdots, x_{k},\cdots)$, where $x_{k} = S_{m,n}^{k}(x) \mod 2$. We can easily see that the following relation of semiconjugacy holds:
	\[
	\pi \circ S_{m,n} = S \circ \pi.
	\]
	As a consequence, for every $S_{m,n}$-invariant measure $\mu$ on $\mathbb{N}$ we can find a $S$-invariant measure $\pi_{*}\mu$ on $\Sigma_{2}$ defined as $\pi_{*}\mu(E) = \mu(\pi^{-1}(E))$.
	
	We observe that for the Syracuse maps $S_{m,n}$ the key topology $\mathcal{T}$ is the discrete one. In fact, we have $\{x\} = \{2x, x\} \cap \{x, S_{m,n}(x)\}$ and we can see that every singleton is intersection of two basic open subsets, then open. As a consequence, every compact subset $C \subset \mathbb{N}$ is finite.
	
	Now let $K \subset \Sigma_{2}$ be the following compact $S$- invariant subset
	\[
	K := \overline{\bigcup \pi(\mathcal{O})} \,\, \text{such that} \,\, \mathcal{O} \,\, \text{is a cycle in} \,\, \mathbb{N}. 
	\]
	We have that the set of accumulation point $K' = \emptyset$ if, and only if, there exist finitely many cycles in $\mathbb{N}$. In fact, once the topology $\mathcal{T}$ in $\mathbb{N}$ is the discrete one, the map $\pi$ is trivially continuous and every set $\pi(\mathcal{O})$ is compact. In the case where we have infinitely many cycles, if $\bigcup \pi(\mathcal{O})$ were compact, then this subset would be closed, which is not true. Hence, we must prove that the boundary $K' = \emptyset$.
	
	Let us suppose, by contradiction that $K' \neq \emptyset$. By Lemma \ref{converse}, since there are infinitely many cycles, there must exist some continuous potential $\phi : \mathbb{N} \to \mathbb{N}_{0}$ with no equilibrium states. 
	
	We define such a potential as follows. Fix $M \in \mathbb{N}$. For a cycle $\mathcal{O}$ with period $p \in \mathbb{N}$ we assign $M$ for $p - 1$ of its points and $0$ for the rest. 
	
	By using the semiconjugacy, we can find a continuous potential $\Phi: \bigcup \pi(\mathcal{O}) \to \mathbb{N}_{0}$ satisfying $\phi = \Phi \circ \pi$. It can be extended to obtain a continuous potential $\Phi_{0}: K \to \mathbb{N}_{0}$. We have $P(\phi) = P(\Phi_{0})$. In fact, by change of variables, we get
	\[
	\int_{\mathcal{O}} \phi d\delta = \int_{\mathcal{O}} \Phi_{0} \circ \pi = \int_{\pi(\mathcal{O})} \Phi_{0} d\pi_{*} \delta.
	\] 
	It implies that
	\[
	P(\Phi_{0}) = \sup_{\delta}\Bigg{\{} \int_{\pi(\mathcal{O})} \Phi_{0} d\pi_{*} \delta\Bigg{\}}.
	\]
	The compactness and invariance of $K$ and continuity of $\Phi_{0}$ guarantees the existence of an equilibrium state $\eta$ whose support is contained in $K'$. Also, we have $\Phi_{0}(\text{supp}\eta)$ compact in $\mathbb{N}_{0}$, then finite. Also $M = \max \Phi_{0}(\text{supp}\eta)$. Since $\{M\}$ is clopen we have the set $U := \Phi_{0}^{-1}(M)$ clopen in $\Sigma_{2}$. By definition of $\phi$ for a cycle $\mathcal{O}_{k}$  with period $p_{k}$ we have $\# \pi(\mathcal{O}_{k}) \bigcap U = p_{k}-1$.
	
	We know that it implies $\int \Phi_{0} \pi_{*} \delta_{k} = M\frac{p_{k} - 1}{p_{k}} \to M$. Hence $P(\Phi_{0}) = M$ and the proportion of point of $\pi(\mathcal{O}_{k})$ inside $U$ converges to $100\%$. The sequence of measures $\pi_{*}\delta_{k}$ has a convergent subsequence $\pi_{*}\delta_{k_{j}}$ whose limit is an equilibrium state for $\Phi_{0}$ and without loss of generality we can assume that it is $\eta$. We then have $\Phi_{0}$ constant equal $M$ in $\text{supp}\eta$ because $\int \Phi_{0} d\eta = M$. 
	
	For any open set $V \supset \text{supp}\eta$ we have infinitely many cycles $\pi(\mathcal{O}_{k_{j}}) \cap V \neq \emptyset$. Otherwise, we do not have $\int \Phi_{0} d\pi_{*}\delta_{k_{j}} \to M$. Taking finitely many cylinders $C_{1}, \cdots, C_{r}$ covering $\text{supp}\eta$ we must have the periodic points of the cycles  $\pi(\mathcal{O}_{k_{j}})$ keeping the initial entries  of the sequence coinciding with the entries of the cylinders up the $(p_{k} - 1)$-th element even when $k \to \infty$. It contradicts the periodicity of the cycles. 
	
	This is a contradiction because we supposed that $K'\neq \emptyset$. It implies that we must have $K' = \emptyset$ and finiteness of cycles in $\mathbb{N}$. This proves the theorem. 
\end{proof}

\end{document}
